% For compatibility with acroread version 
\pdfminorversion=4

\documentclass[mathserif, aspectratio=169]{beamer}
\mode<presentation>

%\usefonttheme[onlymath]{serif}
\usepackage[T1]{fontenc}
%\usepackage{tgheros}
%\usepackage{tgadventor}
%\usepackage{libertine}
%\usepackage{theinhart}
%\usepackage{palatino}
%\usepackage{newtxtext} %Times font only for text

% Use my theme
\usetheme{rognes}

% Misc packages
\usepackage{graphicx}
\usepackage[export]{adjustbox}
\usepackage{stmaryrd}
\usepackage{tikz}
\usepackage{booktabs}
\usepackage{multirow}
\usetikzlibrary{arrows,shapes,backgrounds,decorations,mindmap,patterns,snakes}
\usepackage{amsmath, amssymb, amsthm, amscd}
\usepackage[english]{babel}
\usepackage{algorithm}
\usepackage{pythonhighlight}
\usepackage{epstopdf}
\usepackage{multimedia}
\usepackage{soul}

\newcommand{\R}{\mathbb{R}}
\DeclareMathOperator{\Div}{\mathrm{div}}
\DeclareMathOperator{\Curl}{\mathrm{curl}}
\DeclareMathOperator{\Grad}{\mathrm{grad}}
\DeclareMathOperator{\D}{d}
\newcommand{\foralls}{\forall \;}
\newcommand{\iinner}[2]{\llangle #1, #2 \rrangle}
\newcommand{\inner}[2]{\langle #1,  #2\rangle}
\newcommand{\totald}{\mathrm{d}}
\newcommand{\refer}[1]{\begin{flushright}{\tiny \textcolor{darkgray}{[#1]}}\end{flushright}}
\newcommand{\referinline}[1]{{\tiny \textcolor{darkgray}{[#1]}}}
\newcommand{\ddt}[1]{{#1}_t}
\newcommand{\Mi}{M_i}
\newcommand{\Me}{M_e}
\newcommand{\Iion}{I_{\rm{ion}}}
\newcommand{\dx}{\, \mathrm{d}x}
\newcommand{\ds}{\, \mathrm{d} s}
\newcommand{\triang}{\mathcal{T}}
\DeclareMathOperator{\trace}{tr}
\newcommand{\duality}[2]{\langle #1, #2 \rangle}

\newcommand{\mysection}[1]{\begin{frame} \begin{center} \vspace{3em} \textbf{#1} \end{center} \end{frame}}
\newcommand{\sectionwithrefer}[2]{\begin{frame} \begin{center} \vspace{3em} \textbf{#1} \end{center} \refer{#2} \end{frame}}

\newtheorem{prop}[theorem]{Proposition}


\tikzstyle{na} = [baseline=-.5ex]
\tikzstyle{every picture}+=[remember picture]
\everymath{\displaystyle}

\usepackage[footnotesize, figurename=,size=scriptsize,font={color=rognesblue}]{caption}
%-------------------------------------------------------------------------------
%-------------------------------------------------------------------------------

\DeclareMathOperator{\ssum}{\textstyle \sum}

\begin{document}

%-------------------------------------------------------------------------------
\cleanpage
\begin{frame}
  \begin{columns}
    \begin{column}{0.25\textwidth}
    \centering
    \vspace{1em}
    \includegraphics[width=\textwidth]{graphics/simula_logo_main_RGB.pdf} \\
    \vspace{1em}
    \includegraphics[width=\textwidth]{graphics/waterscales.pdf}  \\
    \vspace{1em}
    \includegraphics[height=0.3\textwidth]{graphics/flag_yellow_low.jpg} 
    \hspace{0.2em}
    \includegraphics[width=0.3\textwidth]{graphics/logo-erc.jpg} \\
    \vspace{1em}
    \includegraphics[width=0.5\textwidth]{graphics/waterscape_logo.png} \\
    \vspace{1em}
    \includegraphics[width=\textwidth]{graphics/rcn-logo.pdf} \\
    \end{column}
    \begin{column}{0.75\textwidth}
    {
    \centering

      \vspace{4em}

      {\bf Brain modelling: \\ from magnetic resonance images \\ to finite element simulation} \\

      \bigskip
      \bigskip
      
      Marie E. Rognes \\
      \bigskip
      \bigskip

      Oslo, Norway \\
      \medskip
      Feb 11 2020 \\
      }
    \end{column}
  \end{columns}
\end{frame}

\cleanpage

\input{/home/meg/presentations/slides/waterscape_uq_tracer_2019.tex}

{\usebackgroundtemplate{\includegraphics[width=\paperwidth]{graphics/Collaborators.pdf}}
\begin{frame}
\end{frame}
}

\begin{frame}
  \frametitle{Outline of lectures}
  \begin{description}
    \item[Lecture 1-2]
      \begin{itemize}
      \item
        Why simulate the brain's waterscape?
      \item
        Where can I find the slides (, book, data, and scripts)?
      \item
        Introduction to brain physiology and imaging
        %% \begin{enumerate}{\roman}
        %% \item
        %%   Introduction to the brain's waterscape: anatomy
        %% \item
        %%   Introduction to the brain's waterscape: mechanics
        %% \item
        %%   Introduction to brain imaging (MRI)
        %% \item
        %%   A software ecosystem: from MRI to finite elements
        %% \end{enumerate}
      \item
        From T1 MRI images to numerical simulation
        %% \begin{itemize}
        %% \item
        %%   A model problem: diffusion of solutes
        %% \item
        %%   From T1 images to a finite element mesh
        %% \item
        %%   A numerical simulation of diffusion
        %% \end{itemize}
      \end{itemize}
    \item[Lecture 3-4]
    \item[Lecture 5-6]
  \end{description}

%\end{enumerate}
\end{frame}

\mysection{Lecture series resources}

\begin{frame}
  \frametitle{Mathematical modeling of the human brain -- book available on GitHub}
  \framesubtitle{Simula SpringerBrief on Computing to be}
  \centering
  \includegraphics[width=\textwidth]{graphics/mri2fem-book-github.png} \\
  \refer{\href{https://github.com/kent-and/mri2fem}{https://github.com/kent-and/mri2fem}}
\end{frame}


\begin{frame}
  \frametitle{Data and software available on Zenodo}
  \vspace{-0.5em}
  \centering
  \includegraphics[width=0.85\textwidth]{graphics/mri2fem-zenodo.png} \\
  \vspace{-1em}
  \refer{\href{https://zenodo.org/communities/mri2fem/}{https://zenodo.org/communities/mri2fem/}}
\end{frame}


\begin{frame}
  \frametitle{Slides and lecture material available on GitHub}
  \centering
  \includegraphics[width=\textwidth]{graphics/mri2fem-lectures-github.png} \\
  \refer{\href{https://github.com/meg-simula/mri2fem-lectures}{https://github.com/meg-simula/mri2fem-lectures}}
\end{frame}

%-------------------------------------------------------------------------------
%
% Lecture 1
% \input{lecture1.tex}
%
%-------------------------------------------------------------------------------

\mysection{The brain's waterscape: anatomy}

\begin{frame}
  \frametitle{Brain anatomy: macro-scale}
\end{frame}

\begin{frame}
  \frametitle{Brain anatomy: meso-scale}
\end{frame}

\begin{frame}
  \frametitle{Brain anatomy: micro-scale}
\end{frame}

\begin{frame}
  \frametitle{Magnetic resonance imaging}
\end{frame}

\mysection{The brain's waterscape: physiology}

\begin{frame}
  \frametitle{Diffusion}
\end{frame}

\begin{frame}
  \frametitle{Glymphatics}
\end{frame}

\begin{frame}
  \frametitle{iPAD}
\end{frame}

\begin{frame}
  \frametitle{Controversies and questions}
\end{frame}


\begin{frame}
  \frametitle{Neurodegenerative disease}
\end{frame}

\begin{frame}
  \frametitle{Sleep}
\end{frame}

\begin{frame}
  \frametitle{Alcohol}
\end{frame}

\begin{frame}
  \frametitle{Exercise}
\end{frame}

\mysection{Introduction to brain magnetic resonance imaging (MRI)}

\begin{frame}
  \frametitle{Magnetic resonance imaging: T1 vs T2}
\end{frame}

\begin{frame}
  \frametitle{Magnetic resonance imaging: DTI}
\end{frame}

\begin{frame}
  \frametitle{Viewing and working with MRI data sets}
\end{frame}

\begin{frame}
  \frametitle{DicomBrowser and example}
\end{frame}

\mysection{A software ecosystem: from MRI to finite elements}

\begin{frame}
\frametitle{Illustration of pipeline}
\end{frame}

\begin{frame}
\frametitle{Packages and installation}
\end{frame}

\begin{frame}
\frametitle{Book Chapter 2, Docker image}
\end{frame}

%-------------------------------------------------------------------------------
%
% Lecture 2
% \input{lecture2.tex}
%
%-------------------------------------------------------------------------------

\mysection{A model problem: diffusion of solutes}

\begin{frame}
\frametitle{}
\end{frame}

\mysection{From T1 images to a finite element mesh}

\begin{frame}
\frametitle{From T1 images to a mesh}

To generate a mesh from an MRI data set including T1-weighted images, we
follow three main steps:
1. extract a T1-weighted image series from the MRI dataset;
2. create (boundary) surfaces from the T1-weighted images using FreeSurfer;
3. generate a volume mesh of the interior using FreeSurfer along with SVM-Tk.
\end{frame}


\begin{frame}
\frametitle{Creating tessellated surfaces}
\end{frame}

\begin{frame}
\frametitle{Refining tessellated surfaces}
\end{frame}

\begin{frame}
\frametitle{Smoothing tessellated surfaces}
\end{frame}

\begin{frame}
\frametitle{Creating volume meshes from surfaces}
\end{frame}

\mysection{A numerical simulation of diffusion}

\begin{frame}
\frametitle{Numerical method}
\end{frame}

\begin{frame}
\frametitle{Set-up}
\end{frame}

\begin{frame}
\frametitle{Numerical results}
\end{frame}


\mysection{Video samples}

%-------------------------------------------------------------------------------


\end{document}
